\subsection{P031 --- Building energy performance assessment using linked data and cross-domain semantic reasoning}
\label{paper:P031}
\textbf{Authors:} Shushan Hu, Jiale Wang, Cathal Hoare, Yehong Li, Pieter Pauwels, James O'Donnell\par
\textbf{Major Category:} BIM and Semantic Information\par
\textbf{Secondary Category:} Building Energy and ZEB\par
\textbf{Keywords:} Building energy performance, Cross-domain semantic reasoning, Semantic web, Data interoperability, Linked data, OWL, RDF\par
\paragraph{主な主張}
RDF/OWLによるlinked dataspaceとcross-domain inference rulesを組み合わせ, sensor, weather, occupancy等にまたがるimplicit knowledgeを推論して, building energy performance metricsに必要な情報を半自動的に構成するframeworkを示した. \cite{Hu:2021CrossDomainSemanticReasoning}
\paragraph{新規性}
BIM由来情報を中心とした既往のsemantic reasoningに対し, meteorological, occupancy, sensing dataを含むcross-domain informationへ推論対象を拡張し, performance assessmentに必要なhidden linksを生成した点に新規性がある.
\paragraph{位置付け}
cross-domain semantic reasoningとimplicit knowledge generationの基礎文献であり, semanticsを予測featureそのものではなく, 評価に必要な情報関係を導出するknowledge layerとして用いる研究として位置付ける.
